\documentclass[../distribution.tex]{subfiles}

\begin{document}
	\noindent Soit $U \subset \mathbb{R}^{n}$ un ouvert bornée tel que $\partial U$ est $\mathscr{C}^{1}$. \\
	\noindent Soit $p \in [1, +\infty]$.

	\vspace{1em}

	L'objectif de cette partir est d'éterdre des fonctions de $W^{1, p}(U)$ à des fonctions de $W^{1, p}(\mathbb{R}^{n})$.

	\mathblock{theorem}{d'extension}{
		Soit $V \subset \subset U$ tel que $U \subset \subset V$. \bigskip \\
		Alors il existe un opérateur linéaire bornée $E$
		$$
			E : W^{1, p}(U) \longrightarrow W^{1, p}(\mathbb{R}^{n})
		$$
		tel que
		\begin{mathlist}
			\item $\forall u \in W^{1, p}(U), (E u)_{| U} = u$ (dans $\mathrm{L}^{p}(U)$).
			\item $\forall u \in W^{1, p}(U), \mathrm{Supp} \ E u \subset V$.
			\item $\exists C > 0, \forall u \in W^{1, p}(U)$,
				$$
					\lVert E u \rVert_{W^{1, p}(\mathbb{R}^{n})} \leq C \lVert u \rVert_{W^{1, p}(U)}
				$$
		\end{mathlist}
	}

	\mathblock{definition}{Extension de $u$ sur $\mathbb{R}^{n}$}{
		Soit $V \subset \subset U$ tel que $U \subset \subset V$. \\
		Soit $E : W^{1, p}(U) \longrightarrow W^{1, p}(\mathbb{R}^{n})$ un opérateur linéaire satisfaisant le théorème précédent. \\
		Soit $u \in W^{1, p}(U)$. \bigskip \\
		On dit que $E u$ est une \textit{extension} de $u$ sur $\mathbb{R}^{n}$.
	}

	Il est également possible d'étendre cette extension pour $k = 2$. On obtiendra alors une majoration telle
	$$
		\lVert E u \rVert_{W^{2, p}(\mathbb{R}^{n})} \leq C \lVert u \rVert_{W^{2, p}(U)}
	$$
\end{document}
